\text{\textbackslash{}documentclass[11pt]\{article\}

\textbackslash{}usepackage[a4paper,margin=1in]\{geometry\}
\textbackslash{}usepackage\{amsmath,amssymb,mathtools\}
\textbackslash{}usepackage\{booktabs\}
\textbackslash{}usepackage\{hyperref\}
\textbackslash{}usepackage\{microtype\}

\textbackslash{}hypersetup\{
  colorlinks=true,
  linkcolor=blue,
  urlcolor=blue
\}

\textbackslash{}title\{\textbackslash{}textbf\{Applying Ding et al.'s Cooling Analysis\}\textbackslash{}\textbackslash{}
to a Long-Range Transverse-Field Ising Model\}
\textbackslash{}author\{\}
\textbackslash{}date\{\}

\textbackslash{}begin\{document\}

\textbackslash{}maketitle

\textbackslash{}section\{The notation problem\}

There are two unrelated quantities that can both be denoted by
\$\textbackslash{}alpha\$:

\textbackslash{}begin\{enumerate\}
    \textbackslash{}item the long-range interaction exponent in the transverse-field
    Ising model;
    \textbackslash{}item the system--bath coupling parameter in the cooling protocols
    of Ding and collaborators.
\textbackslash{}end\{enumerate\}

To avoid ambiguity, this document uses
\$\$
\textbackslash{}alpha\_\{\textbackslash{}mathrm\{LR\}\}
\$\$
for the long-range Ising exponent,
\$\$
\textbackslash{}Gamma
\$\$
for the physical system--bath coupling, and
\$\$
\textbackslash{}alpha\_\{\textbackslash{}mathrm\{SB\}\}
\$\$
for the rescaled coupling used in the follow-up paper.

The relevant papers are:

\textbackslash{}begin\{itemize\}
    \textbackslash{}item \textbackslash{}href\{https://www.alphaxiv.org/abs/2508.05703\}
    \{\textbackslash{}emph\{End-to-End Efficient Quantum Thermal and Ground State Preparation Made Simple\}\},
    by Ding, Zhan, Preskill, and Lin;
    \textbackslash{}item \textbackslash{}href\{https://www.alphaxiv.org/abs/2512.03457\}
    \{\textbackslash{}emph\{Beyond Lindblad Dynamics: Rigorous Guarantees for Thermal and Ground State Preservation under System--Bath Interactions\}\},
    by Wang and Ding.
\textbackslash{}end\{itemize\}

\textbackslash{}section\{Long-range transverse-field Ising model\}

Consider the long-range transverse-field Ising model
\$\$
H\_\{\textbackslash{}mathrm\{sys\}\}(g,\textbackslash{}alpha\_\{\textbackslash{}mathrm\{LR\}\})
=
-J\textbackslash{}sum\_\{i<j\}
\textbackslash{}frac\{\textbackslash{}sigma\_i\textasciicircum{}x\textbackslash{}sigma\_j\textasciicircum{}x\}
\{r\_\{ij\}\textasciicircum{}\{\textbackslash{}alpha\_\{\textbackslash{}mathrm\{LR\}\}\}\}
-h\textbackslash{}sum\_i \textbackslash{}sigma\_i\textasciicircum{}z,
\$\$
where
\$\$
g=\textbackslash{}frac\{h\}\{J\}.
\$\$

The exponent \$\textbackslash{}alpha\_\{\textbackslash{}mathrm\{LR\}\}\$ determines the spatial decay of the
Ising interaction:
\$\$
J\_\{ij\}\textbackslash{}propto r\_\{ij\}\textasciicircum{}\{-\textbackslash{}alpha\_\{\textbackslash{}mathrm\{LR\}\}\}.
\$\$

It is part of the system Hamiltonian. Changing
\$\textbackslash{}alpha\_\{\textbackslash{}mathrm\{LR\}\}\$ changes:

\textbackslash{}begin\{itemize\}
    \textbackslash{}item the many-body energy spectrum;
    \textbackslash{}item the location of the critical point;
    \textbackslash{}item the physical excitation gap;
    \textbackslash{}item the Bohr frequencies;
    \textbackslash{}item the relaxation and mixing properties.
\textbackslash{}end\{itemize\}

For a sweep of \$g=h/J\$, define the physical energy gap
\$\$
\textbackslash{}Delta(g,\textbackslash{}alpha\_\{\textbackslash{}mathrm\{LR\}\})
=
E\_1(g,\textbackslash{}alpha\_\{\textbackslash{}mathrm\{LR\}\})
-
E\_0(g,\textbackslash{}alpha\_\{\textbackslash{}mathrm\{LR\}\}).
\$\$

The minimum gap along the sweep is
\$\$
\textbackslash{}Delta\_\{\textbackslash{}min\}(\textbackslash{}alpha\_\{\textbackslash{}mathrm\{LR\}\})
=
\textbackslash{}min\_\{g\textbackslash{}in[g\_i,g\_f]\}
\textbackslash{}Delta(g,\textbackslash{}alpha\_\{\textbackslash{}mathrm\{LR\}\}).
\$\$

\textbackslash{}section\{The system--bath coupling in Ding et al.\}

In the original Ding--Zhan--Preskill--Lin paper, the total Hamiltonian is
\$\$
H\_\{\textbackslash{}alpha\}(t)
=
H
+
H\_E
+
\textbackslash{}alpha f(t)
\textbackslash{}left(
A\_S\textbackslash{}otimes B\_E
+
A\_S\textasciicircum{}\textbackslash{}dagger\textbackslash{}otimes B\_E\textasciicircum{}\textbackslash{}dagger
\textbackslash{}right).
\$\$

In that equation:

\textbackslash{}begin\{itemize\}
    \textbackslash{}item \$H\$ is the target system Hamiltonian;
    \textbackslash{}item \$H\_E\$ is the bath Hamiltonian;
    \textbackslash{}item \$A\_S\$ is a system coupling operator;
    \textbackslash{}item \$B\_E\$ is a bath operator;
    \textbackslash{}item \$f(t)\$ is a time-dependent interaction envelope;
    \textbackslash{}item \$\textbackslash{}alpha\$ controls the strength of the system--bath interaction.
\textbackslash{}end\{itemize\}

Thus, the \$\textbackslash{}alpha\$ in this paper is not an interaction-range exponent.
It is the coefficient multiplying the system--bath coupling term.

For the single-ancilla bath, the paper uses
\$\$
H\_E=-\textbackslash{}frac\{\textbackslash{}omega\}\{2\}Z,
\textbackslash{}qquad
B\_E=\textbackslash{}frac\{X-iY\}\{2\}
=
|1\textbackslash{}rangle\textbackslash{}langle 0|.
\$\$

The ancilla frequency \$\textbackslash{}omega\$ is sampled so that the bath can couple to
the Bohr frequencies of the system.

\textbackslash{}section\{Notation in the follow-up paper\}

The follow-up paper uses \$\textbackslash{}Gamma\$ instead of \$\textbackslash{}alpha\$ for the physical
system--bath coupling:
\$\$
H\_\{\textbackslash{}Gamma\}(t)
=
H
+
H\_E
+
\textbackslash{}Gamma f(t)
\textbackslash{}left(
A\_S\textbackslash{}otimes B\_E
+
A\_S\textasciicircum{}\textbackslash{}dagger\textbackslash{}otimes B\_E\textasciicircum{}\textbackslash{}dagger
\textbackslash{}right).
\$\$

Here, \$\textbackslash{}Gamma\$ is the physical coupling strength.

The follow-up paper then introduces the rescaled parameter
\$\$
\textbackslash{}alpha\_\{\textbackslash{}mathrm\{SB\}\}
=
\textbackslash{}frac\{\textbackslash{}Gamma\}\{\textbackslash{}sqrt\{\textbackslash{}sigma\}\},
\$\$
where \$\textbackslash{}sigma\$ is the width of the Gaussian filter.

Therefore,
\$\$
\textbackslash{}boxed\{
\textbackslash{}alpha\_\{\textbackslash{}mathrm\{SB\}\}=\textbackslash{}frac\{\textbackslash{}Gamma\}\{\textbackslash{}sqrt\{\textbackslash{}sigma\}\}
\}
\$\$

is a rescaled system--bath parameter. It is unrelated to
\$\textbackslash{}alpha\_\{\textbackslash{}mathrm\{LR\}\}\$.

\textbackslash{}section\{Combined Hamiltonian for the application\}

For the long-range TFIM application, write
\$\$
H\_\{\textbackslash{}mathrm\{sys\}\}(g,\textbackslash{}alpha\_\{\textbackslash{}mathrm\{LR\}\})
=
-J\textbackslash{}sum\_\{i<j\}
\textbackslash{}frac\{\textbackslash{}sigma\_i\textasciicircum{}x\textbackslash{}sigma\_j\textasciicircum{}x\}
\{r\_\{ij\}\textasciicircum{}\{\textbackslash{}alpha\_\{\textbackslash{}mathrm\{LR\}\}\}\}
-Jg\textbackslash{}sum\_i \textbackslash{}sigma\_i\textasciicircum{}z,
\$\$
where \$g=h/J\$.

The total system--bath Hamiltonian is
\$\$
H\_\{\textbackslash{}mathrm\{tot\}\}(t)
=
H\_\{\textbackslash{}mathrm\{sys\}\}(g,\textbackslash{}alpha\_\{\textbackslash{}mathrm\{LR\}\})
+
H\_E
+
\textbackslash{}Gamma f\_\{\textbackslash{}sigma\}(t)
\textbackslash{}left(
A\_S\textbackslash{}otimes B\_E
+
A\_S\textasciicircum{}\textbackslash{}dagger\textbackslash{}otimes B\_E\textasciicircum{}\textbackslash{}dagger
\textbackslash{}right).
\$\$

The notation is then:

\textbackslash{}begin\{center\}
\textbackslash{}begin\{tabular\}\{lll\}
\textbackslash{}toprule
Symbol \& Meaning \& Location \textbackslash{}\textbackslash{}
\textbackslash{}midrule
\$\textbackslash{}alpha\_\{\textbackslash{}mathrm\{LR\}\}\$
\& Long-range Ising exponent
\& \$H\_\{\textbackslash{}mathrm\{sys\}\}\$ \textbackslash{}\textbackslash{}

\$\textbackslash{}Gamma\$
\& Physical system--bath coupling
\& \$H\_\{\textbackslash{}mathrm\{tot\}\}\$ \textbackslash{}\textbackslash{}

\$\textbackslash{}alpha\_\{\textbackslash{}mathrm\{SB\}\}=\textbackslash{}Gamma/\textbackslash{}sqrt\{\textbackslash{}sigma\}\$
\& Rescaled bath coupling
\& Follow-up analysis \textbackslash{}\textbackslash{}

\$\textbackslash{}sigma\$
\& Gaussian filter width
\& \$f\_\{\textbackslash{}sigma\}(t)\$ \textbackslash{}\textbackslash{}

\$g=h/J\$
\& Transverse-field sweep parameter
\& \$H\_\{\textbackslash{}mathrm\{sys\}\}\$ \textbackslash{}\textbackslash{}
\textbackslash{}bottomrule
\textbackslash{}end\{tabular\}
\textbackslash{}end\{center\}

\textbackslash{}section\{Relevant cooling estimates\}

For each value of \$\textbackslash{}alpha\_\{\textbackslash{}mathrm\{LR\}\}\$, calculate both the physical gap
\$\$
\textbackslash{}Delta(g,\textbackslash{}alpha\_\{\textbackslash{}mathrm\{LR\}\})
\$\$
and the dissipative or KMS gap
\$\$
\textbackslash{}lambda\_\{\textbackslash{}mathrm\{KMS\}\}
(g,\textbackslash{}alpha\_\{\textbackslash{}mathrm\{LR\}\},\textbackslash{}sigma).
\$\$

The physical gap controls spectral resolution. The KMS gap controls
dissipative convergence.

For finite-temperature preparation, the beyond-Lindblad analysis gives
the rescaled mixing-time estimate
\$\$
t\_\{\textbackslash{}mathrm\{mix\}\}\textasciicircum{}\{\textbackslash{}mathrm\{rescaled\}\}
\textbackslash{}sim
\textbackslash{}frac\{1\}\{\textbackslash{}lambda\_\{\textbackslash{}mathrm\{KMS\}\}\}
\textbackslash{}log\textbackslash{}left(
\textbackslash{}frac\{\textbackslash{}sigma\textasciicircum{}\{-1/2\}\textbackslash{}beta\}\{\textbackslash{}epsilon\}
\textbackslash{}right),
\$\$
where \$\textbackslash{}beta\$ is the inverse temperature and \$\textbackslash{}epsilon\$ is the target
accuracy.

The discrete channel mixing time is approximately
\$\$
\textbackslash{}tau\_\{\textbackslash{}mathrm\{mix\}\}
\textbackslash{}sim
\textbackslash{}frac\{\textbackslash{}sigma\}\{\textbackslash{}Gamma\textasciicircum{}2\}
t\_\{\textbackslash{}mathrm\{mix\}\}\textasciicircum{}\{\textbackslash{}mathrm\{rescaled\}\}.
\$\$

Using
\$\$
\textbackslash{}alpha\_\{\textbackslash{}mathrm\{SB\}\}
=
\textbackslash{}frac\{\textbackslash{}Gamma\}\{\textbackslash{}sqrt\{\textbackslash{}sigma\}\},
\$\$
this becomes
\$\$
\textbackslash{}tau\_\{\textbackslash{}mathrm\{mix\}\}
\textbackslash{}sim
\textbackslash{}frac\{1\}\{\textbackslash{}alpha\_\{\textbackslash{}mathrm\{SB\}\}\textasciicircum{}2\}
\textbackslash{}frac\{1\}\{\textbackslash{}lambda\_\{\textbackslash{}mathrm\{KMS\}\}\}
\textbackslash{}log\textbackslash{}left(
\textbackslash{}frac\{\textbackslash{}sigma\textasciicircum{}\{-1/2\}\textbackslash{}beta\}\{\textbackslash{}epsilon\}
\textbackslash{}right).
\$\$

Therefore, while the controlled approximation remains valid,
\$\$
\textbackslash{}tau\_\{\textbackslash{}mathrm\{mix\}\}
\textbackslash{}propto
\textbackslash{}alpha\_\{\textbackslash{}mathrm\{SB\}\}\textasciicircum{}\{-2\}.
\$\$

Increasing the bath coupling by a factor of two ideally reduces the
number of channel iterations by a factor of four.

\textbackslash{}section\{Constraint on the bath coupling\}

The perturbative beyond-Lindblad analysis requires higher-order corrections
to remain controlled. The relevant condition is approximately
\$\$
\textbackslash{}alpha\_\{\textbackslash{}mathrm\{SB\}\}\textasciicircum{}2\textbackslash{}sigma
\textbackslash{}lesssim
\textbackslash{}lambda\_\{\textbackslash{}mathrm\{KMS\}\}.
\$\$

Since
\$\$
\textbackslash{}alpha\_\{\textbackslash{}mathrm\{SB\}\}\textasciicircum{}2\textbackslash{}sigma=\textbackslash{}Gamma\textasciicircum{}2,
\$\$
this can also be written as
\$\$
\textbackslash{}Gamma\textasciicircum{}2
\textbackslash{}lesssim
\textbackslash{}lambda\_\{\textbackslash{}mathrm\{KMS\}\}.
\$\$

For the complete sweep, define
\$\$
\textbackslash{}lambda\_\{\textbackslash{}min\}(\textbackslash{}alpha\_\{\textbackslash{}mathrm\{LR\}\})
=
\textbackslash{}min\_\{g\textbackslash{}in[g\_i,g\_f]\}
\textbackslash{}lambda\_\{\textbackslash{}mathrm\{KMS\}\}
(g,\textbackslash{}alpha\_\{\textbackslash{}mathrm\{LR\}\},\textbackslash{}sigma).
\$\$

A conservative coupling bound is
\$\$
\textbackslash{}alpha\_\{\textbackslash{}mathrm\{SB,max\}\}
\textbackslash{}sim
\textbackslash{}sqrt\{
\textbackslash{}frac\{\textbackslash{}lambda\_\{\textbackslash{}min\}(\textbackslash{}alpha\_\{\textbackslash{}mathrm\{LR\}\})\}
\{\textbackslash{}sigma\}
\}.
\$\$

Thus, the long-range exponent affects the maximum usable bath coupling
indirectly through the dissipative gap.

\textbackslash{}section\{Filter width for ground-state preparation\}

For ground-state preparation, the Gaussian filter must resolve the lowest
relevant excitation. The filter width scales approximately as
\$\$
\textbackslash{}sigma
\textbackslash{}gtrsim
\textbackslash{}frac\{1\}\{\textbackslash{}Delta\}
\textbackslash{}log\textbackslash{}left(
\textbackslash{}frac\{\textbackslash{}lVert H\textbackslash{}rVert\}\{\textbackslash{}epsilon\}
\textbackslash{}right),
\$\$
up to logarithmic factors.

For a complete sweep, use the minimum physical gap:
\$\$
\textbackslash{}sigma\_\{\textbackslash{}mathrm\{sweep\}\}
\textbackslash{}gtrsim
\textbackslash{}frac\{1\}\{\textbackslash{}Delta\_\{\textbackslash{}min\}(\textbackslash{}alpha\_\{\textbackslash{}mathrm\{LR\}\})\}
\textbackslash{}log\textbackslash{}left(
\textbackslash{}frac\{\textbackslash{}lVert H\textbackslash{}rVert\_\{\textbackslash{}max\}\}\{\textbackslash{}epsilon\}
\textbackslash{}right).
\$\$

Therefore:

\textbackslash{}begin\{itemize\}
    \textbackslash{}item \$\textbackslash{}alpha\_\{\textbackslash{}mathrm\{LR\}\}\$ affects \$\textbackslash{}sigma\$ through the physical
    minimum gap;
    \textbackslash{}item \$\textbackslash{}alpha\_\{\textbackslash{}mathrm\{SB\}\}\$ affects the number of bath-coupling
    iterations;
    \textbackslash{}item \$\textbackslash{}Gamma\$ is the physical coupling appearing in the Hamiltonian.
\textbackslash{}end\{itemize\}

\textbackslash{}section\{End-to-end cost\}

The total physical evolution time is
\$\$
T\_\{\textbackslash{}mathrm\{total\}\}
=
2T\textbackslash{},\textbackslash{}tau\_\{\textbackslash{}mathrm\{mix\}\},
\$\$
where \$2T\$ is the duration of one system--bath interaction.

Since the Gaussian envelope typically requires
\$\$
T\textbackslash{}gtrsim \textbackslash{}sigma
\$\$
up to logarithmic factors,
\$\$
T\_\{\textbackslash{}mathrm\{total\}\}
\textbackslash{}sim
\textbackslash{}frac\{\textbackslash{}sigma\}\{
\textbackslash{}alpha\_\{\textbackslash{}mathrm\{SB\}\}\textasciicircum{}2
\textbackslash{}lambda\_\{\textbackslash{}mathrm\{KMS\}\}
\}
\textbackslash{}log\textbackslash{}left(
\textbackslash{}frac\{\textbackslash{}sigma\textasciicircum{}\{-1/2\}\textbackslash{}beta\}\{\textbackslash{}epsilon\}
\textbackslash{}right).
\$\$

For a piecewise sweep, divide the path into intervals labeled by \$k\$:
\$\$
g\_k\textbackslash{}longrightarrow g\_\{k+1\}.
\$\$

The cost of interval \$k\$ is approximately
\$\$
T\_k
\textbackslash{}sim
\textbackslash{}frac\{\textbackslash{}sigma\_k\}\{
\textbackslash{}alpha\_\{\textbackslash{}mathrm\{SB\},k\}\textasciicircum{}2
\textbackslash{}lambda\_k
\}
\textbackslash{}log\textbackslash{}left(
\textbackslash{}frac\{\textbackslash{}sigma\_k\textasciicircum{}\{-1/2\}\textbackslash{}beta\}\{\textbackslash{}epsilon\_k\}
\textbackslash{}right),
\$\$
where
\$\$
\textbackslash{}lambda\_k
=
\textbackslash{}lambda\_\{\textbackslash{}mathrm\{KMS\}\}
(g\_k,\textbackslash{}alpha\_\{\textbackslash{}mathrm\{LR\}\},\textbackslash{}sigma\_k).
\$\$

The total sweep cost is
\$\$
T\_\{\textbackslash{}mathrm\{total\}\}\textasciicircum{}\{\textbackslash{}mathrm\{sweep\}\}
\textbackslash{}sim
\textbackslash{}sum\_k T\_k.
\$\$

The segment errors should satisfy
\$\$
\textbackslash{}sum\_k \textbackslash{}epsilon\_k
\textbackslash{}leq
\textbackslash{}epsilon\_\{\textbackslash{}mathrm\{tot\}\}.
\$\$

\textbackslash{}section\{Summary\}

The essential distinction is
\$\$
\textbackslash{}boxed\{
\textbackslash{}text\{\$\textbackslash{}alpha\_\{\textbackslash{}mathrm\{LR\}\}\$ changes the system being cooled,\}
\}
\$\$
whereas
\$\$
\textbackslash{}boxed\{
\textbackslash{}text\{\$\textbackslash{}alpha\_\{\textbackslash{}mathrm\{SB\}\}\$ changes how strongly the bath cools it.\}
\}
\$\$

More explicitly,
\$\$
\textbackslash{}alpha\_\{\textbackslash{}mathrm\{LR\}\}:
\textbackslash{}qquad
J\_\{ij\}\textbackslash{}propto r\_\{ij\}\textasciicircum{}\{-\textbackslash{}alpha\_\{\textbackslash{}mathrm\{LR\}\}\},
\$\$
while
\$\$
\textbackslash{}Gamma:
\textbackslash{}qquad
\textbackslash{}Gamma f(t)
\textbackslash{}left(
A\_S\textbackslash{}otimes B\_E
+
A\_S\textasciicircum{}\textbackslash{}dagger\textbackslash{}otimes B\_E\textasciicircum{}\textbackslash{}dagger
\textbackslash{}right)
\$\$
controls the physical system--bath interaction, and
\$\$
\textbackslash{}alpha\_\{\textbackslash{}mathrm\{SB\}\}
=
\textbackslash{}frac\{\textbackslash{}Gamma\}\{\textbackslash{}sqrt\{\textbackslash{}sigma\}\}
\$\$
is the rescaled parameter used in the follow-up analysis.

The practical workflow is:

\textbackslash{}begin\{enumerate\}
    \textbackslash{}item Choose \$\textbackslash{}alpha\_\{\textbackslash{}mathrm\{LR\}\}\$.
    \textbackslash{}item Sweep \$g=h/J\$ and calculate
    \$\textbackslash{}Delta(g,\textbackslash{}alpha\_\{\textbackslash{}mathrm\{LR\}\})\$.
    \textbackslash{}item Determine \$\textbackslash{}Delta\_\{\textbackslash{}min\}\$ and choose \$\textbackslash{}sigma\$.
    \textbackslash{}item Estimate
    \$\textbackslash{}lambda\_\{\textbackslash{}mathrm\{KMS\}\}
    (g,\textbackslash{}alpha\_\{\textbackslash{}mathrm\{LR\}\},\textbackslash{}sigma)\$.
    \textbackslash{}item Choose \$\textbackslash{}alpha\_\{\textbackslash{}mathrm\{SB\}\}\$ such that
    \$\$
    \textbackslash{}alpha\_\{\textbackslash{}mathrm\{SB\}\}\textasciicircum{}2\textbackslash{}sigma
    \textbackslash{}lesssim
    \textbackslash{}lambda\_\{\textbackslash{}mathrm\{KMS\}\}.
    \$\$
    \textbackslash{}item Estimate
    \$\$
    \textbackslash{}tau\_\{\textbackslash{}mathrm\{mix\}\}
    \textbackslash{}sim
    \textbackslash{}frac\{1\}\{
    \textbackslash{}alpha\_\{\textbackslash{}mathrm\{SB\}\}\textasciicircum{}2
    \textbackslash{}lambda\_\{\textbackslash{}mathrm\{KMS\}\}\}.
    \$\$
    \textbackslash{}item Sum the costs over the \$h/J\$ sweep.
\textbackslash{}end\{enumerate\}

In one sentence:

\$\$
\textbackslash{}boxed\{
\textbackslash{}text\{\$\textbackslash{}alpha\_\{\textbackslash{}mathrm\{LR\}\}\$ changes the Hamiltonian, while
\$\textbackslash{}alpha\_\{\textbackslash{}mathrm\{SB\}\}\$ changes the cooling strength.\}
\}
\$\$

\textbackslash{}end\{document\}}